% 目的
\section{目的}

本プロジェクトの目的は以下の通りである：

\subsection{主目的}

\begin{enumerate}
    \item \textbf{Moduleインターフェース設計パターンの確立} \\
    組み込みシステムにおけるハードウェアモジュール管理の設計パターンを定義し、
    リファレンス実装として公開する。

    \item \textbf{再利用可能なアーキテクチャの提供} \\
    ESP32-S3に限らず、他のマイコンプラットフォームでも適用可能な汎用的な設計として、
    \texttt{IModule}インターフェース、\texttt{ModuleTimer}クラス、\texttt{SystemData}パターンを提供する。
\end{enumerate}

\subsection{副次的目的}

\begin{enumerate}
    \item \textbf{教育・学習用のサンプル} \\
    組み込み開発の初学者が、モジュール化・非ブロッキング処理・データフロー管理を学べる
    具体的なコード例を提供する。

    \item \textbf{ESP32-S3 CAMの活用例} \\
    カメラ、センサー、LED、LCD等の周辺機器を組み合わせた実用的なサンプルアプリケーションを構築する。
\end{enumerate}

\subsection{範囲}

本プロジェクトが提供するもの：
\begin{itemize}
    \item IModuleインターフェースの定義と仕様
    \item ModuleTimerクラスの実装
    \item SystemDataパターンの設計指針
    \item ESP32-S3 CAMを用いたサンプルモジュール群
\end{itemize}
